% ── Chapter 3: Analysis of the Existing System & Methodology ─────────────────
\chapter{Analysis of the Existing System and Methodology}

% =============================================================================
% NOTE — REDACTION STATUS: this chapter is NOT yet redacted.
% The blocks below are the ORGANISED RAW MATERIAL (translated to English) that
% will be turned into prose during redaction. They are kept as % comments so
% nothing half-finished renders in the PDF. Each block sits under the section
% it feeds. Plan/outline lives in guide/02_contribution_guide.md (Chapter 3).
% Convention: chapter Introduction & Conclusion are UNNUMBERED (\section*).
% =============================================================================


% ─────────────────────────────────────────────────────────────────────────────
\section*{Introduction}
\addcontentsline{toc}{section}{Introduction}

Part~I set out the foundations of this work: how organisations turn data into decisions, and why
logistics --- data-rich and thin-margined --- is a domain in which Business Intelligence delivers
particular value. This chapter opens Part~II by grounding those principles in a concrete setting.
Before any warehouse can be modelled or any dashboard built, the existing situation must first be
understood: who the actors are, which systems already run, how the business operates day to day, and
precisely where the analytical gap lies.

The chapter therefore proceeds as a structured analysis. It introduces the organisations behind the
project, studies the systems and processes already in place, and from that diagnosis derives the
requirements, the prioritised business scope, and the methodology that guide the remainder of the
contribution:

\begin{itemize}[noitemsep]
    \item \textbf{Host and client organisations} (\S\ref{sec:ch3-orgs}) --- Ourquilane, the software
          company carrying the project, and Yalidine, the logistics operator that is its final client,
          together with its national network.
    \item \textbf{Study of the existing systems} (\S\ref{sec:ch3-existing}) --- the services and
          business processes in operation, the source systems that record them, and the analytical gap
          they leave behind.
    \item \textbf{Requirements specification} (\S\ref{sec:ch3-requirements}) --- the functional and
          non-functional requirements derived from that gap.
    \item \textbf{Business scope and prioritisation} (\S\ref{sec:ch3-scope}) --- the two business axes,
          ranked with the MoSCoW method.
    \item \textbf{Development methodology and planning} (\S\ref{sec:ch3-methodology}) --- the iterative
          approach adopted and the project timeline.
\end{itemize}


% ═════════════════════════════════════════════════════════════════════════════
\section{Host and Client Organisations}
\label{sec:ch3-orgs}

% -----------------------------------------------------------------------------
\subsection{Ourquilane — The Project Carrier}
% NOTES FOR REDACTION (Ourquilane):
%   - A dynamic startup specialised in developing innovative software solutions,
%     based in Hydra, Algiers.
%   - Mission: transform and modernise various sectors through digital platforms,
%     digital marketing, and decision-support tools.
%   - Product portfolio (gives credibility + shows the logistics-IT focus):
%       * HR Force   — HR management (used by Yalidine)
%       * FLEET GO   — fleet management
%       * ITOP       — IT management
%       * CASH Box   — financial operations
%       * Route Planner — route / tour planning
%       * CS Care    — complaints management
%   - Ourquilane is the entity CARRYING this project (la societe porteuse).
%   - [VISUAL-PLACEHOLDER] Ourquilane organisation chart (organigramme) + logo.

% -----------------------------------------------------------------------------
\subsection{Yalidine El Djazair Services — The Client}
% NOTES FOR REDACTION (Yalidine — the final client of the project):
%   - Leader of express delivery in Algeria; provides cutting-edge national logistics.
%   - Core services: B2B courier, B2C e-commerce delivery, and cash-on-delivery
%     collection (COD = Cash On Delivery).
%   - Recognised for real-time tracking ensuring the security and speed of shipments.
%   - Certifications: ISO 9001 (Quality), ISO 14001 (Environment),
%     ISO 27001 (Information Security).
%   - [VISUAL-PLACEHOLDER] Yalidine logo (Figure "Logo de l'organisme client").
%
% PRINCIPAL SERVICES OF YALIDINE  --> these become the project's TWO BUSINESS AXES:
%   (1) Express parcel delivery (B2C / e-commerce) — last-mile delivery - Pricipal service.
%   (2) On-demand dedicated transport (B2B) — companies transporting goods/parcels/Papers.
%   NB: introduce them here as Yalidine's real services; the MoSCoW prioritisation
%       (§3.4) is what frames them as Must/Should-have axes.

% -----------------------------------------------------------------------------
\subsubsection{The Logistics Network}
% NOTES FOR REDACTION (Yalidine network / operational structure):
%
% Organisational layer:
%   - General direction + 3 regional directions
%   - Call centres
%   - Vehicle fleet
%
% Three-tier operational structure (Hub -> Regional centre -> Station):
%   * HUB (Main Sorting Centre / Centre de Tri Principal):
%       The principal platforms of the national network. They receive parcels
%       from various origins, perform major sorting, then redistribute them to
%       the Regional Centres. Yalidine has FOUR strategically located hubs:
%         - Oued Smar (Algiers) : covers the north of the country
%         - Relizane            : covers the west
%         - Constantine         : covers the east
%         - Ghardaia            : covers the Great South
%   * REGIONAL SORTING CENTRE (Centre de tri regional):
%       A logistics entity that manages several delivery stations. Receives
%       consolidated parcels from the hubs and organises their distribution to
%       its attached stations.
%   * STATION (Centre / unite de proximite):
%       The proximity unit responsible for LAST-MILE delivery. Each station
%       covers a set of communes and is the unique attachment point of its
%       drivers. Its main mission: receive parcels from the regional centre and
%       organise the precise assignment of parcels to drivers for the daily
%       tours. This is where volume-forecast precision matters most.
%       Yalidine has a dense network of 182 stations covering the 55 wilayas.
%
% Delivery ZONES (drive pricing/distance):
%       - Zone 0 : intra-wilaya
%       - Zone 1 : inter-wilaya, short distance
%       - Zone 2 : inter-wilaya, long distance
%       - Zone 3 : inter-wilaya, very long distance (south)
%
% Delivery TYPES:
%       - SD = Stop Desk (collection at a desk/point)
%       - HD = Home Delivery
%
% KEY FIGURES (decide whether to include; CONFIRM before printing):
%   % TODO(achraf): confirm Yalidine figures before they appear in the PDF
%       - 1 general direction
%       - 3 regional directions
%       - 3000+ employees
%       - 1000+ vehicles
%       - 4 hubs
%       - 180+ delivery stations
%
%   - [VISUAL-IMPLEMENT] TikZ diagram of the Hub -> Regional Centre -> Station
%     hierarchy (optionally annotated with the 4 hubs and their coverage).
%   - [VISUAL-IMPLEMENT] grid table: delivery zones (0-3) and delivery types (SD/HD).


% ═════════════════════════════════════════════════════════════════════════════
\section{Study of the Existing Systems}
\label{sec:ch3-existing}

% -----------------------------------------------------------------------------
\subsection{Business Processes}
% NOTES FOR REDACTION (the existing processes — ONE subsection, organised per service):
%
% A) PARCEL DELIVERY (B2C e-commerce):
%   - Parcel lifecycle at Yalidine (cycle de vie d'un colis): from creation/pickup,
%     through hub & regional sorting, to station and last-mile delivery (or return).
%   - Parcel statuses (status d'un colis): the multiple states a parcel moves
%     through across its lifecycle.  % TODO(achraf): list the exact status set
%   - Parcel pricing / tariffs (tarification d'un colis):
%       * delivered-parcel tariff (tarif colis livre)
%       * return-parcel tariff (tarif colis retour)
%       * theoretical tariff vs real tariff  (theoretical = grid; real = f(volume, weight))
%       * parcel reimbursement (remboursement — tied to COD in case of assurance pay 1% of the value of the parcel, reimbursed in case damaged or lost, lost 100% , bu damaged depends on the embalage status of the parcel from 60% to 80%).
%   - [VISUAL-IMPLEMENT] TikZ parcel-lifecycle / status flow diagram.
%
% B) DEDICATED TRANSPORT (B2B, on-demand):
%   - Process: transport request -> Pricing and agremment on pricing -> assigned dedicated trip -> pickup -> delivery -> completion.
%   % TODO(achraf): confirm the exact transport request -> stops workflow (you can develop it yourself)

% -----------------------------------------------------------------------------
\subsection{Source Systems}
% NOTES FOR REDACTION (existing IT systems of Yalidine):
%   - Present each at a FUNCTIONAL level (what it provides), not internals.
%   - Describe them as the REAL source systems (the simulation is introduced
%     only in Ch.6 — do NOT mention simulation here).
%
% The FIVE internal systems the warehouse consumes:
%   1. Core logistics system  (Yalidine App)
%        -> geography, delivery centres, pricing grid, full parcel-event history
%   2. HR system              (HR Force)
%        -> companies, agencies, employees, organisational hierarchy
%   3. Cash-box system        (CashBox)
%        -> operating expenses, freelance-driver payments, reimbursements, transfers
%   4. Payroll system         (PC Paie)
%        -> monthly payslips
%   5. Transport system
%        -> dedicated B2B transport requests and their stops
%
%   ( Returly App -> returns management — exists at Yalidine but is NOT part of
%     the consumed source set; mention only if useful, otherwise omit. )
%
%   - [VISUAL-IMPLEMENT] grid table: source system -> domain it provides
%     (5 rows, columns: System | Domain provided).

% -----------------------------------------------------------------------------
\subsection{Identified Weaknesses and the Analytical Gap}
% NOTES FOR REDACTION:
%   - Data is scattered across the disconnected source systems above; no single
%     consolidated, analytical view.
%   - What decision-makers CANNOT do today: get a timely, reliable, cross-system
%     view of each axis along Operations / Cost & Profitability / Performance,
%     without slow, manual extraction and reconciliation.
%   - This analytical gap is the justification for the BI platform (ties back to
%     the General Introduction problem statement).


% ═════════════════════════════════════════════════════════════════════════════
\section{Requirements Specification}
\label{sec:ch3-requirements}

% -----------------------------------------------------------------------------
\subsection{Functional Requirements}
% NOTES FOR REDACTION (short, numbered list derived from the gap):
%   - TWO axis pages, each with THREE sub-treatments:
%       Operations / Cost & Profitability / Performance     (NO "Pricing" page)
%   - Plus: Overview page, Alerts page, Administration page.
%   - RBAC user model: users IMPORTED from HR Force, ACTIVATED by a super-admin,
%     each assigned a role that controls dashboard access.
%   - Proactive alerting system, tied to the KPIs (threshold rules -> notifications).
%   - [VISUAL-IMPLEMENT] use-case diagram (TikZ)  OR  [VISUAL-PLACEHOLDER] drawn UML.

% -----------------------------------------------------------------------------
\subsection{Non-Functional Requirements}
% NOTES FOR REDACTION (qualities the platform must meet):
%   - Performance (responsive analytics over large volumes — aggregate layer).
%   - Scalability (production-scale data: hundreds of millions of parcel events).
%   - Security (stateless token auth, RBAC).
%   - Usability (clear decision-oriented dashboards, KPI cards).
%   - Maintainability / evolvability (modular, layered architecture).
%   - Cost: open-source stack, no licensing cost.


% ═════════════════════════════════════════════════════════════════════════════
\section{Business Scope and Prioritisation (MoSCoW)}
\label{sec:ch3-scope}
% NOTES FOR REDACTION:
%   - Must have   -> On-Demand Dedicated Transport (B2B; Dedicated Trip / Courier / Handling)
%   - Should have -> Parcel Delivery (classic e-commerce express)
%   - Could have  -> Route Analysis  (ONE neutral line: future work; data not
%                    available within project scope — no GPS/confidentiality talk)
%   - Won't have  -> (state explicitly what is out of scope for this version)
%   - Two-independent-financial-perimeters principle: the costs and revenues of
%     the two axes are NEVER mixed (each axis read on its own perimeter).
%   - [VISUAL-IMPLEMENT] MoSCoW grid table: priority -> axis -> what it delivers.


% ═════════════════════════════════════════════════════════════════════════════
\section{Development Methodology and Planning}
\label{sec:ch3-methodology}
% NOTES FOR REDACTION:
%   - Development approach: iterative.
%   - Project timeline / schedule from thesis.md (the PFE planning).
%   - Tools used (link forward to the Ch.4 technology stack at a high level).
%   - [VISUAL-IMPLEMENT] Gantt-style timeline (pgfgantt) — DECIDE: keep initial +
%     final Gantt, or fall back to a milestone table if a Gantt is not kept.
%     % TODO(achraf): decide Gantt vs milestone table; needs pgfgantt package if Gantt.


% ─────────────────────────────────────────────────────────────────────────────
\section*{Conclusion}
\addcontentsline{toc}{section}{Conclusion}
% TODO(redaction): summarise the analysis (organisations, existing systems, the
% gap, the prioritised scope, the methodology) and transition to Chapter 4
% (Design & Architecture).


% =============================================================================
% PARKED REFERENCE (added to references.bib as bourzam2025 — cite where relevant,
% e.g. when discussing demand forecasting / volume prediction at station level):
%   Bourzam & Belhadj, "Prediction des Demandes de Livraison via Apprentissage
%   Profond et Detection du Concept Drift", Memoire de fin d'etudes, ESI Alger.
% =============================================================================
